\section{何を「実施できる」とみなしたか}
\label{sec:conditions}

打ち上げを決めるのは当日の朝である。そこで、毎朝 9 時（JST）に次の 4 つを順に確かめ、
すべて満たした日を「実施できる日」と数えた。1 日 1 回、9 時に上げる想定である。

\begin{enumerate}
  \item \textbf{打ち上げ地点に雪がない}。ERA5 の積雪深（水当量）が \SI{1}{\milli\metre} 未満。
    新雪ならおよそ \SI{1}{\centi\metre} に当たる。
  \item \textbf{地上の風と雨が打ち上げに向く}。9 時の地上 \SI{10}{\metre} の風速が
    \SI{5}{\metre\per\second} 以下で、9--12 時の降水が \SI{0.5}{\milli\metre} 未満。
    風が強いと気球に水素・ヘリウムを詰める作業ができず、放球の瞬間に機体が地面を打つ。
  \item \textbf{落下予測が目標に入る}。目標は次のどれかである。
    \begin{itemize}
      \item 平野: 標高 \SI{100}{\metre} 以下・傾斜 \ang{2} 以下の陸地で、人口集中地区（DID）の外
      \item 瀬戸内海: 関門・豊予（関崎〜佐田岬）・鳴門・紀淡の各海峡の内側
      \item 海岸から \SI{3}{\kilo\metre} 以内の海（「海岸から数 km」の中間の値。\ref{sec:sensitivity} 節で 0--10 km に変えた）
    \end{itemize}
  \item \textbf{陸の落下点に雪がない}。判定は 1. と同じ。
\end{enumerate}

3. で落下予測が山地や沖の海に入る日は、打ち上げを見送る日として数える。
実際にも前日から当日朝の予報で落下点を計算し、目標に入らなければ見送るはずだからである。

打ち上げ場所そのものについては、「電柱やビルがない開けた場所（グラウンドのある公園は可）」
という条件がある。これは約 \SI{5}{\kilo\metre} 間隔の格子では判定できないので、
候補地を絞ったあとに OpenStreetMap で周辺の広場を探した（\ref{sec:osm} 節）。
格子の段階では、打ち上げ候補を「陸地で、標高 \SI{200}{\metre} 以下」に限った。
山の上の開けた場所は、雪と車での到達の点で冬には使いにくいためである。

\begin{whybox}{飛行条件を 2 通り比べた理由}
落下点までの距離は飛行時間でほぼ決まる。標準の飛行（上昇 \SI{5}{\metre\per\second}・
破裂 \SI{30}{\kilo\metre}・降下は海面で \SI{6}{\metre\per\second}）は約 2.3 時間飛び、
冬は東へ約 \SI{255}{\kilo\metre} 流される。短時間の飛行（上昇 \SI{8}{\metre\per\second}・
破裂 \SI{20}{\kilo\metre}）は約 1.3 時間で、流される距離は約 \SI{180}{\kilo\metre} に縮む。
どちらが有利かは、打ち上げ地点から東にどんな地形が並んでいるかで変わる。
\end{whybox}
